
Combining formal verification strategies for LLM code generation appears straightforward---stack grammar constraints, static analysis, and SMT-guided repair for multiplicative error reduction. However, this work reveals a critical prerequisite overlooked by prior approaches: while error detection is model-agnostic, repair via feedback loops is model-capability-dependent. A completion model given static analysis feedback increased security issues by 600\%, not decreased them.

This finding matters for practitioners deploying verification pipelines in production. Teams investing in multi-stage verification infrastructure may observe code quality degradation if model capability is not matched to each stage's requirements.

\subsubsection{1.1 The Problem}

At the surface level, the problem is well-documented: LLM-generated code contains syntax errors, security vulnerabilities, and specification violations. Different verification strategies---grammar-constrained decoding, static analysis, SMT-guided repair---target different error types. Each strategy has been studied in isolation: grammar constraints reduce compilation errors by over 50\% (Mündler et al., 2025), static analysis feedback reduces security issues from over 40\% to 13\% \citep{blyth2025static}, and SMT-guided pipelines achieve 75--82\% pass@1 on contract synthesis \citep{lim2025contracteval}.

However, no unified comparison exists across all three strategies on identical benchmarks with consistent metrics. This leaves a critical gap: it is unknown whether error classes are truly independent (enabling additive or multiplicative gains) or overlapping (yielding diminishing returns). Without this knowledge, practitioners cannot predict combined pipeline effectiveness.

\subsubsection{1.2 Contributions}

This work makes the following contributions:

\begin{enumerate}
\item \textbf{Error-class-independence framework.} The first quantitative measurement of independence between grammar, semantic, and specification error classes, using Jaccard indices on improvement sets across 164 HumanEval problems. All pairwise Jaccard indices fall below the 0.30 threshold, with a mean of 0.0752.
\end{enumerate}

\begin{enumerate}
\item \textbf{Model capability prerequisite identification.} Identification that feedback-based verification stages require instruction-tuned models; completion models may degrade output quality. This distinction---detection is model-agnostic, repair is model-capability-dependent---was not previously documented.
\end{enumerate}

\begin{enumerate}
\item \textbf{Partial pipeline validation.} Demonstration that grammar-constrained decoding reduces syntax errors by 40 percentage points on a proof-of-concept set, validating the first stage of the layered verification pipeline.
\end{enumerate}
